% oja_template.tex
% Unofficial LaTeX template for publishing 
% in the Open Journal of Astrophysics
% v1.0 released September 6, 2015 (matches openjournal.cls)
% Author: Emmanuel Frion

% Basic setup
\documentclass[twocolumn,apj,numberedappendix]{openjournal}
% Available options:
% [twocolumn] - two-column mode
% [onecolumn] - (default) main text in one-column mode
% [apj]       - typeset in the style of ApJ.
% [apjl]      - (default) typeset in the style of ApJ Letters 
% [tighten]   - some adjustments to approximate grid typesetting
% [numberedappendix]   - number appendix sections as A, B, etc
% [appendixfloats]  - use separate numbering for floats within appendix
% [twocolappendix]  - make appendix in two-col mode in a two-col paper
% [revtex4]   - force using revtex4 (rather than 4-1)

% In case of issues with doi formatting, uncomment these lines
% \let\olddoi\doi
% \renewcommand{\doi}[1]{\href{https://doi.org/#1}{DOI: \nolinkurl{#1}}}

% Optional useful packages
\usepackage[dvipsnames]{xcolor}
% \usepackage{textgreek}
\usepackage[T1]{fontenc}
\usepackage[utf8]{inputenc}
\usepackage[english]{babel}
\usepackage{natbib}
\usepackage{amsmath}
\usepackage{listings}
\definecolor{lbcolor}{rgb}{0.9,0.9,0.9}
\definecolor{cosmiclatte}{rgb}{1.0, 0.97, 0.91}
\lstset{language=Python,
        basicstyle=\footnotesize\ttfamily,
        showspaces=false,
        showstringspaces=false,
        tabsize=2,
        breaklines=false,
        breakatwhitespace=true,
        backgroundcolor = \color{cosmiclatte},
        identifierstyle=\ttfamily,
        keywordstyle=\bfseries\color[rgb]{0.133,0.545,0.133},
        commentstyle=\color[rgb]{0.133,0.545,0.133},
        stringstyle=\color[rgb]{0.627,0.126,0.941},
        frame=single
    }   
    
\usepackage{hyperref}
\hypersetup{
    unicode, 
    colorlinks=true,
    linkcolor=linkcolor,
    citecolor=linkcolor,
    filecolor=linkcolor,
    urlcolor=linkcolor,
}
\usepackage{color,colortbl}
\definecolor{linkcolor}{rgb}{0.0,0.3,0.5}
\DeclareGraphicsExtensions{.bmp,.png,.jpg,.pdf}
\usepackage{orcidlink}

\usepackage{savesym}
\savesymbol{tablenum}
\usepackage{siunitx}
\restoresymbol{SIX}{tablenum}
\usepackage{float}
\usepackage{marginnote}
\usepackage{tabularray}
\usepackage{booktabs}
\usepackage{bm}
\usepackage{tikz}

% this should be loaded last
\usepackage{microtype}

\UseTblrLibrary{booktabs}
\usetikzlibrary{calc}

\urlstyle{same}

\sisetup{mode=match, per-mode=symbol, uncertainty-mode=separate}
\DeclareSIUnit{\erg}{erg}

\input{newcommands.inc}

\floatplacement{figure}{tp}
\floatplacement{table}{tp}

\journalinfo{The Open Journal of Astrophysics}

\begin{document}

\title{Nuclear Reaction Integration on GPUs}

\author{Michael Zingale$^1$\orcidlink{0000-0001-8401-030X}}

\affiliation{$^1$Department of Physics and Astronomy,
             Stony Brook University,
             Stony Brook, NY 11794-3800, USA}


\begin{abstract}
We discuss the challenges on performing nuclear reaction integration
on GPUs, considering different integrators, precision, and applications.
\end{abstract}

\begin{keywords}
    {nucleosynthesis}
\end{keywords}

\maketitle


\section{Introduction}

% introduction

Discuss previous integration studies

New challenges posed by GPUs: vectorization / branching, register
pressure, trend toward reduced precision.

\section{Integrators}


Backward-Euler and BDF / VODE

Rosenbrock (YASS is a special case)

Jacobian structure / sparsity

Jacobian matters more with Rosenbrock.


\section{Tests}

\subsection{One-zone burn}

\subsection{Reduced precision}

\subsection{Flame}

\subsection{Detonation}

look at number of RHS evals vs. time

\subsection{Double-detonation SN Ia}


\section{Summary and Future Directions}



\begin{acknowledgments}
The work at Stony Brook was supported by
DOE/Office of Nuclear Physics grant DE-FG02-87ER40317.
\end{acknowledgments}

\software{AMReX \citep{amrex,amrex-gpu},
          AMReX-Astro Microphysics \citep{Microphysics,microphysics-joss},
          autodiff \citep{autodiff},
          Jupyter \citep{jupyter},
          matplotlib \citep{Hunter:2007},
          NetworkX \citep{networkx},
          numba \citep{numba},
          numpy \citep{numpy,numpy2020},
          pynucastro \citep{pynucastro,pynucastro2},
          pytest \citep{pytest},
          SciPy \citep{scipy},
          SymPy \citep{sympy},
          VODE \citep{vode},
          yt \citep{yt}}

\bibliographystyle{aasjournal}

\bibliography{refs}

\end{document}
